\subsection{Torsion-character flags by finite flat cotangent sequences}
\label{sec:geometry-torsion-character-flags}
The following local calculation separates the finite flat algebra from
the geometric inputs to the test-point theorem of
Section~\ref{sec:geometry-relative-chern-comparison}.
It applies to the original coordinate flag, not just a flag adapted to
a valuation-ring basis.

Let \(R\) be the valuation ring of an algebraically closed complete
rank-one mixed-characteristic field \(C\), with additive valuation
\(v(p)=1\). Put \(F=\mathbf F_{p^f}\), \(q=p^f\), and
\(N=F^n\), and let \(G\) be a finite locally free commutative
\(F\)-module scheme over \(R\), with its specified generic
identification \(G(C)=N\). Write \(M=\mathcal O(G)\subset R^N\)
for its evaluation lattice. We use the fundamental Teichm\"uller
characters \(\chi_\tau:F\to R\), extended by zero at zero, indexed
by the Frobenius orbit, so
\(\chi_{\phi\tau}=\chi_\tau^p\).
Write \(F_i\subset N\) for the first \(i\) coordinate subspace and
\(G_i\subset G\) for its schematic closure (\(0\le i\le n\)).
These closures are finite locally free: their coordinate algebras
are the finitely generated torsion-free images of \(M\) under
restriction to \(F_i\). Over a valuation ring such modules are free.
The group and field actions descend because the corresponding
generic identities hold in torsion-free algebras.
We use the classical finite locally free subgroup quotient theorem
(for commutative finite locally free groups), so \(G/G_i\) is
finite locally free and \(G\to G/G_i\) is a \(G_i\)-torsor.
In particular these assertions do not require a discrete valuation.

Let \(A_\tau\subset R^n\) consist of coefficient vectors \(a\)
such that \(\sum_j a_j\chi_\tau(x_j)\) is in \(M+pR^N\).
Let \(D_{\tau i}\) be its coordinate-\(i\) image on vectors with
coordinates \(<i\) zero. Let \(I_{\tau i}\subset R\) consist of
\(c\) such that \(c\chi_\tau(x_i)\) on \(F_{i+1}\) is the
restriction of a function of \(M\). These are the original
\texttt{af}, \texttt{dId} and \texttt{xId}, respectively.
All these coefficient lattices are finitely generated: the
conditions are images or kernels of finite matrices over a
valuation ring, for which finite Smith elimination applies.
All \(D_{\tau i}\) and \(I_{\tau i}\) are nonzero principal ideals.
For the first statement below, \(w\in R\) satisfies
\(w^{p-1}\in pR^\times\).

\paragraph{1. Annihilator and discriminant.}
One has \(wR^n\subset A_\tau\) for every \(\tau\).
In fact, for each \(\mathbf F_p\)-linear form
\(\mu:N\to\mathbf F_p\), the generic character
\(\zeta_p^{\mu(x)}\) extends to \(G\): it is a generic point of
the finite Cartier dual, and \(R\) is integrally closed.
Set
\[
 \ell(z)=\sum_{a=1}^{p-1}(-1)^{a+1}(z-1)^a/a,\qquad
 u=\zeta_p-1 .
\]
The formal logarithm identity through degree \(p-1\), whose
coefficients are \(p\)-integral, gives
\(\ell(\zeta_p^b)\equiv b\ell(\zeta_p)\pmod{pR}\)
for \(0\le b<p\): terms of degree at least \(p\) in \(u\)
have valuation \(>1\). Moreover \(v(\ell(\zeta_p))=1/(p-1)\).
Thus \(w\mu\) has a lift modulo \(p\) in \(M\), using integer
representatives of \(\mu\). In an \(\mathbf F_p\)-basis of \(F\),
each \(\chi_\tau(x_j)\) is a linear combination of such
coordinate functions modulo \(p\), since its lift is additive
modulo \(p\). This proves the inclusion. Evaluation at the
coordinate vectors proves \(A_\tau\subset R^n\).

For any of the groups \(P\) in this discussion, write
\(\lambda(P)=v(\det\mathcal O(P))\) in the full evaluation
lattice on its split generic points, and let
\(d(P)=\operatorname{length}_v\omega_{P/R}\).
Here the latter length of a finitely presented torsion module
is the sum of valuations of its nonzero Smith invariants.
There is an exact identity
\begin{equation}\label{eq:torsion-character-discriminant}
  2\lambda(P)=|P(C)|d(P).
\end{equation}
To see this without any tightness hypothesis, first recall why
cotangent lengths are additive. Finite locally free group schemes
are local complete intersections (SGA 3, Exp.~III, 4.15).
Their cotangent complexes at the identity are two-term complexes
of finite locally free modules in degrees \(-1,0\).
In the generically \'{e}tale case the degree-\(-1\)
differential is injective: its generic localization is an
isomorphism and its domain is torsion-free. In an exact finite
locally free group sequence \(H\to P\to Q\), transitivity for
the cotangent complex, and flat base change to the identity
fibre \(H\) of \(P\to Q\), give
\[
 0\longrightarrow\omega_Q\longrightarrow\omega_P
       \longrightarrow\omega_H\longrightarrow0 .
\]
All maps are the original differential maps. In particular
\(d(P)=d(H)+d(Q)\), also after applying any unit-group
character idempotent.

For a rank-one \(F\)-module scheme \(Q_i=G_{i+1}/G_i\),
the nonconstant augmentation character components all have rank
one. By Raynaud's classification (Cor.~1.5.1 of
\emph{Sch\'emas en groupes de type \((p,\ldots,p)\)},
Bull.~SMF 102 (1974), 241--280), if
\(I_{\tau i}=c_{\tau i}R\), the fundamental generators
\(X_\tau=c_{\tau i}\chi_\tau\) give the presentation
\[
 \mathcal O(Q_i)=
 R[X_\tau:\tau]/(X_\tau^p-\delta_{\tau i}X_{\phi\tau}:\tau),
 \qquad
 \delta_{\tau i}=c_{\tau i}^{p}/c_{\phi\tau,i}.
\]
The coefficient ideals here are indeed the specified
\(I_{\tau i}\): a function on \(G_{i+1}\) constant on the
generic \(G_i\)-cosets is invariant under \(G_i\) on the
whole product, by flatness. Faithfully flat descent identifies
it with a function on \(Q_i\). Conversely every function on
\(Q_i\) pulls back to \(G_{i+1}\) and lifts to \(M\).

The presentation implies, with the indicated coefficient
identification and not just up to isomorphism of lengths,
\begin{equation}\label{eq:torsion-character-cotangent}
 \omega_{Q_i,\tau}=I_{\tau i}/I_{\phi^{-1}\tau,i}^{\otimes p}.
\end{equation}
Here the tensor power is its product ideal in \(R\).
The monomials with exponents \(0,\ldots,p-1\) in the
\(X_\tau\) form a basis. Their unscaled evaluation matrix
is a unit matrix in determinant, because the nonconstant
monomials enumerate the \(q-1\) multiplicative characters
on \(F^\times\), extended by zero at zero. Consequently
\[
 \lambda(Q_i)=\tfrac12 q(p-1)\sum_\tau v(c_{\tau i}),
 \qquad d(Q_i)=(p-1)\sum_\tau v(c_{\tau i}).
\]
For a finite locally free \(H\)-torsor \(P\to Q\), the relative
discriminant ideal is the extension of the discriminant of
\(H/R\): this can be checked after the faithfully flat
base change trivializing the torsor. Transitivity of the
trace discriminant then makes
\(\lambda(P)/|P(C)|\) additive in \(H,Q\).
Indeed the generic trace pairing on the split coordinate
algebra is the dot product, whose Gram determinant in an
integral evaluation basis has valuation \(2\lambda(P)\).
Induction on the schematic coordinate flag, using the
rank-one formula, proves
\eqref{eq:torsion-character-discriminant} for \(G\) and every
coordinate quotient \(G/G_i\).

\paragraph{2. The two exact step-graph identities.}
We specify the still-required geometric interface to the
original graph \(\Theta_\tau\).
Assume its first projection is the augmentation ideal \(M^+\),
its second projection is \(L_\tau^\vee\), and its second-coordinate
kernel is \(pL_\tau^\vee\), as in \texttt{Base.th}.
Assume also the three original chart pins: the graph contains
\(((M^+)^2,0)\); it is equivariant under \(F^\times\) acting by
\(g(x)\mapsto g(ax)\) and by multiplication by \(\chi_\tau(a)\)
on the second coordinate; and on the \(\chi_\tau\)-eigenspace
\((g,0)\) in the graph implies \(g\in(M^+)^2\).
These are exactly the pins of
\texttt{chartD\_th\_sqIdeal\_mem\_and\_smul\_mem\_and\_eig\_mem\_sqIdeal}.
They imply the actual map of the graph modulo \(p\) is
\[
 M^+\longrightarrow\omega_{G,\tau}
               \xrightarrow{\ \sim\ } L_\tau^\vee/pL_\tau^\vee.
\]
Indeed the kernel over a fixed first coordinate is
\(pL_\tau^\vee\); apply the character projector (with the
unit denominator \(q-1\)) to obtain surjectivity and
injectivity on the indicated augmentation quotient.

Let \(K_{\tau i}\subset L_\tau^\vee\) be the second
projection of the graph restricted to functions vanishing on
\(F_i\). Since restriction \(M\to\mathcal O(G_i)\) is onto,
the image modulo \(p\) is precisely the kernel of
\(\omega_{G,\tau}\to\omega_{G_i,\tau}\). This follows directly
by lifting a sum of products in the augmentation square.
The preceding exact cotangent sequence gives
\[
 K_{\tau i}/K_{\tau,i+1}
       \simeq\omega_{Q_i,\tau}
       =I_{\tau i}/I_{\phi^{-1}\tau,i}^{\otimes p}.
\]
For the original step graph \(B_{\tau i}\), the map to this
quotient sends \(c\) to the differential class of
\(c\chi_\tau(x_i)\) restricted to \(G_{i+1}\).
Its first projection is \(I_{\tau i}\); for any lift of this
restricted function to \(M\) the class of its second coordinate
modulo \(K_{\tau,i+1}\) is independent of the lift.
Conversely all elements of \(K_{\tau,i+1}\) can be added to the
second coordinate using functions zero on \(F_{i+1}\).
The preceding quotient isomorphism therefore proves the
literal identities
\[
 \operatorname{pr}_2 B_{\tau i}=K_{\tau i},\qquad
 \ker(\operatorname{pr}_2:B_{\phi\tau,i}\to K_{\phi\tau,i})
     =I_{\tau i}^{\otimes p}.
\]
The last kernel is identified via its first coordinate.
This establishes two of the three identities of RC2 without
any product equality of cotangent lengths of the full order
or its intermediate flags. It does not replace the chart pins
by an arbitrary abstract isomorphism of cotangent modules.

\paragraph{3. The coordinate ideal equality from quotient generation.}
Call a function affine modulo \(p\) if, after subtraction of
its value at zero, it is additive on the split points modulo
\(pR\). Suppose the coordinate algebra of \(G/G_i\) is generated
over \(R\) by such functions, for every \(0\le i<n\). This is
precisely the quotient-generation premise \texttt{GenAF}:
invariant functions descend by the finite flat torsor, and
conversely pull back. Under this single premise we prove
\begin{equation}\label{eq:torsion-character-coordinate}
 D_{\tau i}=I_{\tau i}
\end{equation}
without a further full-order tightness hypothesis.

Let \(J_{\tau i}\subset I_{\tau i}\) consist of coefficients
of restrictions to \(Q_i\subset G/G_i\) of affine-modulo-\(p\)
functions in the augmentation \(\chi_\tau\)-eigenspace of
\(\mathcal O(G/G_i)\). Pullback to \(G\) shows
\(J_{\tau i}\subset D_{\tau i}\): such a function modulo
\(p\) is a linear combination of the \(\chi_\tau(x_j)\) with
\(j\ge i\). For clarity the fundamental characters give an
\(R/pR\)-basis of additive coordinate functions, since their
evaluation matrix on an \(\mathbf F_p\)-basis of \(F\)
is invertible modulo the maximal ideal. Its reduction is a
Moore matrix: a nonzero additive polynomial of degree at
most \(p^{f-1}\) cannot vanish at all \(p^f\) field elements.
Projection onto the chosen character is permitted since
\(q-1\) is a unit. The coefficient at \(i\) can be taken to
be the actual value on the \(i\)-th coordinate vector.

By generation and restriction to the closed subgroup \(Q_i\),
the differentials of these functions generate its cotangent
module. Subtract their values at zero and project onto the
\(\chi_\tau\)-component; projection preserves additivity
modulo \(p\). Formula \eqref{eq:torsion-character-cotangent}
gives
\[
 I_{\tau i}=J_{\tau i}+I_{\phi^{-1}\tau,i}^{\otimes p}.
\]
Whenever the second ideal is strictly smaller than the first,
this forces \(J_{\tau i}=I_{\tau i}\), since the first ideal is
principal over a local valuation ring. At an equality step
the result propagates along Frobenius: the \(p\)-th power of
an affine-modulo-\(p\) augmentation eigenfunction is again
one, at the successor character. If every step is an equality
then all \(I_{\tau i}=R\), by cycling the valuations. In that
case \(Q_i\) is constant by the presentation. At least one
\(J_{\tau i}\) must be the unit ideal: otherwise all the
generating functions would restrict to constants modulo the
maximal ideal on this constant subgroup. Indeed after removal
of constants their reductions are additive and hence sums
of their fundamental projections. This would contradict
generation of its full residue function algebra. Propagating
once around the orbit now proves \(J_{\tau i}=I_{\tau i}\)
in every case. We obtain \(I_{\tau i}\subset D_{\tau i}\).

For the reverse conclusion in degrees, let \(m=fn\),
\(h=|N|=p^m\), and let \(A\subset R^m\) be the full additive
coefficient lattice, expressed in any \(\mathbf F_p\)-coordinate
basis. It contains \(wR^m\), is contained in \(R^m\), and
decomposes into the \(A_\tau\) by the unit character projectors
and the invertible Moore change of basis above. Put
\[
 \alpha=v(\det A)=\sum_{\tau,i}v(D_{\tau i}).
\]
In a Smith basis \(\epsilon_j b_j\) for \(A\), the \(b_j\)
are the rows of a unit matrix, and
\(0\le v(\epsilon_j)\le 1/(p-1)\).
Choose lifts \(g_j\in M\) of these linear forms modulo \(p\).
For \(p>2\), the normalized functions \(g_j/\epsilon_j\)
are integral and reduce to the independent linear forms
\(b_j\). For \(p=2\) the same argument works unless
\(v(\epsilon_j)=1\); in this case take the exact lift by
linear combination of the functions \(2\mu(x)=1-(-1)^{\mu(x)}\),
which belong to \(M\) by Cartier duality.

The evaluation matrix of the monomials in the normalized
functions, with each exponent \(0,\ldots,p-1\), is invertible
over \(R\). In the residue field this follows by the invertible
linear change from the \(m\) original coordinates: the new
linear forms generate the full function algebra, and their
\(p\)-th powers are linear combinations of themselves, so the
\(p^m\) indicated monomials span and form a basis.
The unnormalized monomials belong to \(M\) and have
evaluation determinant valuation
\(\frac12 h(p-1)\alpha\).
It follows that
\[
 \lambda(G)\le\tfrac12 h(p-1)\alpha,\qquad
 \alpha\ge d(G)/(p-1)=\sum_{\tau,i}v(I_{\tau i})
\]
by \eqref{eq:torsion-character-discriminant} and cotangent
additivity. The inclusions already proved give the opposite
inequality term by term, and hence equality in every term.
Principal ideals in a valuation ring with the same valuation
are equal. This proves \eqref{eq:torsion-character-coordinate}.

The same determinant calculation gives a convenient sufficient
condition for the quotient-generation premise itself. For any
coordinate quotient \(P\), if its additive coefficient lattice
satisfies \((p-1)\alpha(P)=d(P)\), the above unnormalized
monomials have exactly the determinant valuation of
\(\mathcal O(P)\); their inclusion is therefore equality.
They generate the order by affine-modulo-\(p\) functions.
Thus just the original additive-coefficient degree equality
on these quotient orders suffices, without separate
discriminant identities. This last equality is not asserted
for arbitrary finite flat groups by the argument.

In summary, the annihilator inclusion is automatic for the
finite flat order, the two step-graph identities use only the
literal chart pins, and the remaining coordinate comparison
uses quotient generation (or its sufficient determinant
equality). The proof replaces the full \texttt{FlagConcl}
length-and-ideal comparison as an input to the relative
Chern calculation. At the original geometric test points one
must still realize the specified order, the chart pins and
quotient generation, and the exact Hodge image and kernel
comparisons RC1. In particular the existence of an induced
finite flat group alone does not prove all \texttt{TestFacts}.
